\documentclass[12pt]{article}

\usepackage{amsmath, amssymb}
\usepackage{geometry}
\usepackage{booktabs}
\usepackage{hyperref}

\geometry{margin=1in}

\title{CSE400 -- Fundamentals of Probability in Computing}
\author{Lecture Scribe: Joint Random Variables and Joint Distributions}
\date{Lectures L15--L18 \\ Instructor: Dr. Dhaval Patel}

\begin{document}

\maketitle

\section{Motivation for Joint Random Variables}

In many practical systems, more than one random variable must be analyzed simultaneously. Studying a single random variable is often insufficient when the behavior of one variable depends on another.

The lecture introduces \textbf{joint random variables} to model such relationships.

\section{Examples Motivating Joint Random Variables}

Several applications discussed in the lecture illustrate the need for joint analysis.

\subsection{Smart Transportation}

Random variables:

\begin{itemize}
\item $X =$ Vehicle Speed
\item $Y =$ Traffic Density
\end{itemize}

Question discussed:

Can vehicle speed be high when traffic density is also high?

Observation:  
Speed and traffic density can be \textbf{correlated random variables}.

Applications mentioned in the lecture:

\begin{itemize}
\item Adaptive traffic signals
\item Congestion prediction
\item Autonomous vehicle routing
\end{itemize}

\subsection{Wireless Network Performance}

Random variables:

\begin{itemize}
\item $X =$ Signal-to-Noise Ratio (SNR)
\item $Y =$ Data Throughput
\end{itemize}

Throughput depends on SNR but may also be affected by network conditions such as congestion.

Question discussed:

If SNR is high, will throughput always be high?

Applications:

\begin{itemize}
\item 5G scheduling
\item Adaptive modulation
\item Network planning
\end{itemize}

\subsection{Weather Prediction}

Random variables:

\begin{itemize}
\item $X =$ Rainfall
\item $Y =$ Temperature
\end{itemize}

Application: weather forecasting.

\subsection{Drone Delivery}

Random variables:

\begin{itemize}
\item $X =$ Wind Speed
\item $Y =$ Delivery Time
\end{itemize}

Application: UAV logistics systems.

\subsection{Dice Experiment}

Experiment: roll two dice.

Define:

\begin{itemize}
\item $X =$ outcome of the first die
\item $Y =$ outcome of the second die
\end{itemize}

Questions such as

\[
X + Y = 7
\]

or

\[
X > Y
\]

require analysis using a \textbf{joint probability distribution}.

\section{Joint Random Variables}

When two random variables $X$ and $Y$ are considered together, their probabilities are described by a \textbf{joint distribution}.

This distribution specifies the probability that both variables take specific values simultaneously.

\section{Discrete Case: Joint Probability Mass Function (Joint PMF)}

For discrete random variables $X$ and $Y$, the joint distribution is defined by the \textbf{joint probability mass function}.

\subsection*{Definition}

\[
p_{X,Y}(x,y) = P(X=x, Y=y)
\]

This represents the probability that

\begin{itemize}
\item $X = x$
\item $Y = y$
\end{itemize}

simultaneously.

\subsection*{Properties}

\begin{enumerate}
\item \textbf{Non-negativity}

\[
p_{X,Y}(x,y) \ge 0
\]

\item \textbf{Total probability}

\[
\sum_x \sum_y p_{X,Y}(x,y) = 1
\]

\end{enumerate}

\section{Joint PMF Table Representation}

The joint PMF for discrete variables can be represented using a \textbf{two-dimensional table}.

\begin{center}
\begin{tabular}{c|ccc}
$Y \backslash X$ & $x_1$ & $x_2$ & $x_3$ \\
\hline
$y_1$ & $p(x_1,y_1)$ & $p(x_2,y_1)$ & $p(x_3,y_1)$ \\
$y_2$ & $p(x_1,y_2)$ & $p(x_2,y_2)$ & $p(x_3,y_2)$ \\
$y_3$ & $p(x_1,y_3)$ & $p(x_2,y_3)$ & $p(x_3,y_3)$
\end{tabular}
\end{center}

Each entry corresponds to

\[
P(X=x_i, Y=y_j)
\]

\section{Example: Rolling Two Dice}

Consider the experiment of rolling two dice.

Define:

\begin{itemize}
\item $X =$ result of die 1
\item $Y =$ result of die 2
\end{itemize}

Possible values:

\[
X,Y \in \{1,2,3,4,5,6\}
\]

Total number of outcomes:

\[
6 \times 6 = 36
\]

Each outcome can be represented as an ordered pair:

\[
(x,y)
\]

Examples:

\[
(1,1), (1,2), (1,3), \ldots, (6,6)
\]

If both dice are fair:

\[
P(X=x, Y=y) = \frac{1}{36}
\]

for every valid pair.

The joint PMF therefore forms a \textbf{$6 \times 6$ table} of probabilities.

\section{Continuous Case: Joint Probability Density Function (Joint PDF)}

For continuous random variables, the joint distribution is described using a \textbf{joint probability density function}.

Let

\[
f_{X,Y}(x,y)
\]

denote the joint PDF.

\subsection*{Probability over a Region}

The probability that the pair $(X,Y)$ lies within a region $R$ is

\[
P((X,Y) \in R) =
\int\int_R f_{X,Y}(x,y)\,dx\,dy
\]

\subsection*{Properties}

\begin{enumerate}
\item \textbf{Non-negativity}

\[
f_{X,Y}(x,y) \ge 0
\]

\item \textbf{Total probability}

\[
\int_{-\infty}^{\infty}\int_{-\infty}^{\infty} f_{X,Y}(x,y)\,dx\,dy = 1
\]

\end{enumerate}

\section{Geometric Interpretation}

For continuous random variables, the joint PDF can be viewed as a surface over the $(x,y)$ plane.

The probability that $(X,Y)$ lies within a region corresponds to the \textbf{volume under the density surface above that region}.

\section{Joint Cumulative Distribution Function (Joint CDF)}

The \textbf{joint cumulative distribution function} describes cumulative probabilities for two random variables.

\subsection*{Definition}

\[
F_{X,Y}(x,y) = P(X \le x, Y \le y)
\]

This represents the probability that both events occur simultaneously:

\begin{itemize}
\item $X \le x$
\item $Y \le y$
\end{itemize}

\section{Joint CDF -- Continuous Case}

For continuous random variables with joint PDF $f_{X,Y}(x,y)$,

\[
F_{X,Y}(x,y)
=
\int_{-\infty}^{x}
\int_{-\infty}^{y}
f_{X,Y}(u,v)\,dv\,du
\]

\section{Joint CDF -- Discrete Case}

For discrete variables,

\[
F_{X,Y}(x,y)
=
\sum_{u \le x} \sum_{v \le y}
P(X=u, Y=v)
\]

\section{Properties of the Joint CDF}

\subsection*{1. Bounds}

\[
0 \le F_{X,Y}(x,y) \le 1
\]

\subsection*{2. Monotonicity}

The function is non-decreasing in each variable.

If

\[
x_1 \le x_2
\]

then

\[
F_{X,Y}(x_1,y) \le F_{X,Y}(x_2,y)
\]

Similarly for $y$.

\subsection*{3. Limit Property}

As both variables approach infinity,

\[
F_{X,Y}(x,y) \rightarrow 1
\]

\subsection*{4. Boundary Property}

If either variable approaches negative infinity, the probability approaches zero.

Example:

\[
F_{X,Y}(-\infty, y) = 0
\]

\end{document}